% Native semantic reconstruction of source physical page 26 (printed page 25).
% Provenance: IEEE Std 1516-2025, Clause 3.1, definitions region template--scaffolding object class description.
\HLAStartSemanticPage{26}{25}
\setlength{\HLASemanticDefinitionGap}{12.10pt}
\fontsize{10.02pt}{11.524pt}\selectfont
\begingroup
\hbadness=10000\relax
\addfontfeatures{FakeStretch=1.035}
\enlargethispage{12pt}
\vspace*{2.95pt}

\begin{hladefinition}[0pt]{def-region-template}{region template}
A region template is an incomplete region specification where one or more dimensions have not been assigned ranges.
\end{hladefinition}

\begin{hladefinition}{def-register}{register}
To invoke the \HLAService{Register Object Instance} or the \HLAService{Register Object Instance With Regions} service to create a unique object instance designator.
\HLADefinitionNote[11.58pt]{See 6.8 of IEEE Std 1516.1\HLAStandardMark-2025.}
\end{hladefinition}

\begin{hladefinition}{def-registered-class}{registered class}
The object class that was an argument to the Register Object Instance or the Register Object Instance With Regions service invocation that resulted in the creation of the object instance designator for a given object instance.
\end{hladefinition}

\begin{hladefinition}{def-regular-interaction-class-description}{regular interaction class description}
A description of an interaction class that follows\linebreak IEEE Std 1516.2\HLAStandardMark-2025 and contains at a minimum the interaction class name and publish/subscribe\linebreak indicator.
\end{hladefinition}

\begin{hladefinition}{def-regular-object-class-description}{regular object class description}
A description of an object class that follows IEEE Std 1516.2\HLAStandardMark-2025 and contains at a minimum the object class name and publish/subscribe indicator.
\end{hladefinition}

\begin{hladefinition}{def-remove}{remove}
To receive an invocation of the \HLAService{Remove Object Instance} \HLAPrintersDagger{} service for a particular object instance.
\HLADefinitionNote[11.58pt]{See 6.15 of IEEE Std 1516.1\HLAStandardMark-2025.}
\end{hladefinition}

\begin{hladefinition}{def-repeated-description}{repeated description}
A description of a construct (e.g., object class, interaction class, data type, dimension) in a federation object model (FOM) module with a name and additional properties that are identical to a description that is already available in the current FOM.
\end{hladefinition}

\begin{hladefinition}{def-resolution}{resolution}
The smallest resolvable value separating attribute or parameter values that can be discriminated. Resolution may vary with magnitude for certain data types.
\end{hladefinition}

\begin{hladefinition}{def-retraction}{retraction}
An action performed by a federate to unschedule a previously scheduled message. Message retraction may be visible to the federate to whom the scheduled message was to be delivered. Retraction is widely used in classic event-oriented discrete event simulations to model behaviors such as preemption and interrupts.
\end{hladefinition}

\begin{hladefinition}{def-runtime-infrastructure-rti}{runtime infrastructure (RTI)}
The software that provides common interface services during a High Level Architecture (HLA) federation execution for synchronization and data exchange.
\end{hladefinition}

\begin{hladefinition}{def-runtime-infrastructure-rti-initialization-data-rid}{runtime infrastructure (RTI) initialization data (RID)}
RTI vendor-specific information needed to run an RTI. If required, a RID is supplied when an RTI is initialized.
\end{hladefinition}

\begin{hladefinition}{def-scaffolding-interaction-class-description}{scaffolding interaction class description}
A description of an interaction class in a federation object model (FOM) module or simulation object model (SOM) module that follows IEEE Std 1516.2\HLAStandardMark-2025. However, it will contain only the name. The name of the scaffolding interaction class description is identical to the name of a regular interaction class description provided by another FOM/SOM module. Scaffolding interaction class descriptions act as placeholders in order to represent the class hierarchy from the other FOM/SOM module(s) from which a new regular interaction class description can be specified.
\end{hladefinition}

\begin{hladefinition}{def-scaffolding-object-class-description}{scaffolding object class description}
A description of an object class in a federation object model (FOM) module or simulation object model (SOM) module that follows IEEE Std 1516.2\HLAStandardMark-2025. However, it will contain only the name. The name of the scaffolding object class description is identical to the name of a regular object class description provided by another FOM/SOM module. Scaffolding object class descriptions act as placeholders in order to represent the class hierarchy from the other FOM/SOM module(s) from which a new regular object class description can be specified.
\end{hladefinition}
\endgroup
\clearpage
